\documentclass[12pt]{article}
\usepackage{hyperref}
\usepackage{amsmath, amssymb}
\usepackage{geometry}
\geometry{margin=1in}
\title{Modern Review: Bayesian vs Frequentist Statistics}
\author{ ChatGPT }
\date{\today}

\begin{document}

\maketitle

\section{Core Philosophical Differences}

\textbf{Frequentist:}
\begin{itemize}
    \item Probability is defined as long-run frequency of repeatable experiments.
    \item Parameters are fixed but unknown; probability applies only to data.
    \item Inference uses estimates, \textit{p-values}, and confidence intervals.
    \item Hypothesis tests evaluate whether data are unlikely under a null model.
\end{itemize}

\textbf{Bayesian:}
\begin{itemize}
    \item Probability is a degree of belief about hypotheses and parameters.
    \item Parameters are random variables with prior distributions; data updates beliefs via Bayes' theorem.
    \item Produces posterior distributions, direct probabilities for hypotheses, and credible intervals.
\end{itemize}

\section{Practical Differences in Modern Use}

\subsection{Interpretation}
\begin{itemize}
    \item Bayesian credible interval: 95\% probability the parameter lies in the interval.
    \item Frequentist confidence interval: in repeated experiments, 95\% of intervals contain the true parameter.
\end{itemize}

\subsection{Hypothesis Testing}
\begin{itemize}
    \item Frequentist: relies on p-values and reject/accept decisions.
    \item Bayesian: uses Bayes factors or posterior odds to compare models/hypotheses.
\end{itemize}

\section{Computational and Modeling Trends}

\subsection{Computation}
\begin{itemize}
    \item Bayesian methods: historically intensive; modern MCMC, HMC, and Variational Inference make complex models practical.
    \item Frequentist methods: often analytic and computationally efficient.
\end{itemize}

\subsection{Model Flexibility}
\begin{itemize}
    \item Bayesian hierarchical models express complex dependency structures and multi-level uncertainty.
    \item Frequentist models can handle complexity but may require heavier assumptions.
\end{itemize}

\section{Strengths and Limitations}

\begin{tabular}{|l|l|l|}
\hline
Aspect & Bayesian & Frequentist \\
\hline
Interpretation of uncertainty & Direct probability statements & Long-run frequency interpretation \\
Use of prior knowledge & Explicit & Does not use prior info \\
Adaptive/sequential analysis & Natural updating & Requires careful design \\
Small sample performance & Often superior with good priors & Can perform poorly \\
Computational cost & Intensive for complex models & Efficient \\
Community adoption & Growing in ML, AI, adaptive trials & Dominant in traditional stats/regulatory \\
\hline
\end{tabular}

\section{Modern Trends and Hybrid Approaches}
\begin{itemize}
    \item Empirical Bayes: frequentist estimation of prior hyperparameters, then Bayesian inference.
    \item Bayesian optimization in machine learning for hyperparameter tuning.
    \item Bayesian methods provide uncertainty estimates; frequentist models give point estimates.
\end{itemize}

\section{When to Choose Which}

\textbf{Bayesian preferred when:}
\begin{itemize}
    \item Prior information exists.
    \item Sequential updating is needed.
    \item Rich uncertainty quantification required.
    \item Complex hierarchical/latent structures.
\end{itemize}

\textbf{Frequentist preferred when:}
\begin{itemize}
    \item Objective tests without subjective input.
    \item Computing simplicity and tradition matter.
    \item Regulatory or classical scientific standards require standard tests.
\end{itemize}

\section{Key Textbooks and Resources}

\subsection{Comparison-Focused}
\begin{itemize}
    \item \textit{A Comparison of the Bayesian and Frequentist Approaches to Estimation} --- Francisco J. Samaniego
    \item \textit{Bayesians Versus Frequentists: A Philosophical Debate on Statistical Reasoning} --- Jordi Vallverd\'u
    \item \textit{Bayesian and Frequentist Regression Methods} --- Jon Wakefield
\end{itemize}

\subsection{Core Bayesian Texts}
\begin{itemize}
    \item \textit{Statistical Rethinking} --- Richard McElreath
    \item \textit{Bayesian Statistical Methods} --- Brian J. Reich \& Sujit K. Ghosh
\end{itemize}

\subsection{Classic References}
\begin{itemize}
    \item \textit{Bayesian Data Analysis} --- Andrew Gelman et al.
    \item \textit{Doing Bayesian Data Analysis} --- John K. Kruschke
    \item \textit{All of Statistics} --- Larry Wasserman
\end{itemize}

\subsection{Free PDFs and Open Resources}
\begin{itemize}
    \item Bayesian Data Analysis PDF (Columbia site) --- \url{https://sites.stat.columbia.edu/gelman/book/}
    \item OpenIntro Statistics (free PDF) --- \url{https://www.openintro.org}
    \item Free Bayesian lecture notes and tutorials online.
\end{itemize}

\end{document}
